% ════════════════════════════════════════════════════════════════
\section{Wykład 4 i 5}
\subsection{Programowanie dynamiczne}
% ════════════════════════════════════════════════════════════════

Czasami gdy zaimplementujemy jakąś formę rekurencyjną dostajemy algorytm działający w czasie wykładniczym. 
Powodem może być to, że jakieś wartości obliczane są wielokrotnie.

\paragraph{Problem liczb Fibonacciego}
\begin{flalign*}
	\begin{cases}
		F_0 = F_1 = 1 & n = 0, 1 \\[0.5ex]
		F_n = F_{n - 1} + F_{n - 2} & n \geq 2
	\end{cases} && \text{Zadanie: Obliczyć } n\text{-tą liczbę Fibonacciego} &&
\end{flalign*}

\begin{alg}{Algorytm naiwny Fib1}{alg:fib1}
	\FunctionNN{Fib1}{$n$}
	\If{$n \leq 1$} \Return 1 
	\EndIf
	\State \Return \Call{Fib1}{$n - 1$} $+$ \Call{Fib1}{$n - 2$}
	\EndFunction
\end{alg}

Zobaczmy jak jest obliczane np. \Call{Fib1}{$5$}. W drzewie wywołań poszczególne wartości powtarzają się:
\begin{flalign*}
	& \begin{array}{lc}
		F(i) \text{ jest obliczane} & \text{razy} \\
		(1) & 5 \\
		(2) & 3 \\
		(3) & 2 \\
		(4) & 1 \\
		(5) & 1
	\end{array} && &&
\end{flalign*}

Niech $s(n)$ będzie liczbą dodawań wykonanych przez alg. w wywołaniu \Call{Fib1}{$n$}.
\begin{flalign*}
	& \begin{cases}
		s(0) = s(1) = 0 \\
		s(n) = s(n - 1) + s(n - 2) + 1 & n \geq 2
	\end{cases} && && \\[1.5ex]
	& \text{Niech } f(n) = s(n) + 1 && && \\[1.5ex]
	& \text{Wtedy } f(0) = f(1) = 1 && && \\[1.5ex]
	& f(n) = s(n) + 1 = s(n - 1) + s(n - 2) + 1 + 1 = f(n - 1) - 1 + f(n - 2) - 1 + 1 + 1 && && \\
	& \phantom{f(n)} = f(n - 1) + f(n - 2) && && \\[1.5ex]
	& \text{Zatem } f(n) = F_n \qquad f(n) = \BT(F_n) = \BT\left(\left(\frac{1 + \sqrt{5}}{2}\right)^n\right) = \BM(1.61^n). \text{ Zatem } s(n) = \BM(1.61^n). && &&
\end{flalign*}
\Call{Fib1}{} jest algorytmem wykładniczym.

Oto alg. progr. dynamicznego:
\begin{alg}{Algorytm dynamiczny Fib2}{alg:fib2}
	\FunctionNN{Fib2}{$n$} \Comment{$F$ - tablica}
	\If{$n \leq 1$} \Return 1 
	\EndIf
	\State $F(0) \Assign F(1) \Assign 1$
	\For{$i \Assign 2$ \To $n$}
	\State $F(i) \Assign F(i - 1) + F(i - 2)$
	\EndFor
	\State \Return $F(n)$
	\EndFunction
\end{alg}

Liczba dodawań wykonanych przez \Call{Fib2}{} w wywołaniu \Call{Fib2}{$n$} wynosi $n - 1$.
Złożoność czasowa wynosi $\BT(n)$.

% ════════════════════════════════════════════════════════════════
\subsection{Algorytm mnożenia macierzy}
% ════════════════════════════════════════════════════════════════

\begin{alg}{Naiwne mnożenie macierzy}{alg:matmul}
	\FunctionNN{MatrixMultiply}{$A, B$}
	\Comment{$A$ -- macierz $p \times q$; $B$ -- macierz $q \times r$; $C$ -- macierz $p \times r$}
	\For{$i \Assign 1$ \To $p$}
	\For{$j \Assign 1$ \To $r$}
	\State $C(i, j) \Assign 0$
	\For{$k \Assign 1$ \To $q$}
	\State $C(i, j) \Assign C(i, j) + A(i, k) \cdot B(k, j)$ \Comment{operacja dominująca}
	\EndFor
	\EndFor
	\EndFor
	\State \Return $C$
	\EndFunction
\end{alg}

Operacją dominującą jest mnożenie w linii 5. Mamy $p \cdot q \cdot r$ mnożeń. Złożoność jest $\BO(p \cdot q \cdot r)$.

% ════════════════════════════════════════════════════════════════
\subsection{Problem mnożenia ciągu macierzy}
% ════════════════════════════════════════════════════════════════

\begin{example}{Mnożenie ciągu macierzy}
$(A_1, A_2, A_3)$ gdzie $A_1$ -- macierz $5 \times 10$, $A_2$ -- macierz $10 \times 3$, $A_3$ -- macierz $3 \times 2$.
Ile potrzeba mnożeń skalarnych (mnożenia liczb) do obliczenia iloczynu $A_1 \cdot A_2 \cdot A_3$?
\begin{flalign*}
	A_1 \cdot A_2 \cdot A_3 = (A_1 \cdot A_2) \cdot A_3 = A_1 \cdot (A_2 \cdot A_3) && &&
\end{flalign*}

\begin{itemize}
	\item[I sposób:] $(A_1 \cdot A_2) \cdot A_3 = 5 \cdot 10 \cdot 3 + 5 \cdot 3 \cdot 2 = 150 + 30 = 180$
	\item[II sposób:] $A_1 \cdot (A_2 \cdot A_3) = 10 \cdot 3 \cdot 2 + 5 \cdot 10 \cdot 2 = 60 + 100 = 160$
\end{itemize}
II sposób jest szybszy.
\end{example}

\paragraph{Problem:}
Dla danego ciągu macierzy $(A_1, A_2, \dots, A_n)$, gdzie $A_i$ jest macierzą $p_{i-1} \times p_i$, znajdź kolejność wykonywania mnożeń minimalizującą liczbę mnożeń skalarnych użytych do obliczenia iloczynu $A_1 \cdot A_2 \cdot \dots \cdot A_n$.

Ile jest różnych sposobów obliczenia tego iloczynu? Innymi słowy, na ile sposobów możemy poprawnie wstawić nawiasy w iloczynie $A_1 \cdot A_2 \cdot \dots \cdot A_n$?

\begin{example}{Kolejności nawiasowania dla $n=4$}
$A_1 \cdot A_2 \cdot A_3 \cdot A_4$ \quad ($n = 4$)
\begin{itemize}
	\item $((A_1 \cdot A_2) \cdot A_3) \cdot A_4$
	\item $(A_1 \cdot (A_2 \cdot A_3)) \cdot A_4$
	\item $A_1 \cdot ((A_2 \cdot A_3) \cdot A_4)$
	\item $A_1 \cdot (A_2 \cdot (A_3 \cdot A_4))$
	\item $(A_1 \cdot A_2) \cdot (A_3 \cdot A_4)$
\end{itemize}
Zatem mamy 5 sposobów.
\end{example}

Sprawdźmy czy umiemy rozwiązać ten problem szybko poprzez sprawdzenie wszystkich możliwości.
Niech $P(n)$ -- liczba sposobów wstawienia nawiasów w iloczynie $n$ macierzy.
Załóżmy, że ostatnie mnożenie w ciągu $A_1 \dots A_n$ to mnożenie między $A_k$ i $A_{k+1}$:
$(A_1 \dots A_k)(A_{k+1} \dots A_n)$. W podciągach lewym i prawym nawiasy są rozmieszczane niezależnie od siebie.

Stąd rekurencja:
\begin{flalign*}
	& \begin{cases}
		P(n) = \sum_{k=1}^{n-1} P(k) \cdot P(n - k) \\
		P(1) = 1
	\end{cases} && &&
\end{flalign*}
Jest to znana formuła rekurencyjna a rozwiązanie to przesunięta liczba Catalana:
\begin{flalign*}
	P(n) = C(n - 1), \text{ gdzie } C(n) = \frac{1}{n + 1}\binom{2n}{n} \text{ jest n-tą liczbą Catalana.} && &&
\end{flalign*}
Wiadomo, że $C(n) = \BM\left(\frac{4^n}{n^{3/2}}\right)$. Zatem $P(n)$ jest wykładniczy względem $n$.
Chcemy znaleźć szybki algorytm.

\paragraph{Rozwiązanie z optymalną podstrukturą:}
Zauważmy że jeśli optymalne rozwiązanie jest postaci $(A_1 \dots A_k)(A_{k+1} \dots A_n)$ dla jakiegoś $k \in \{1, \dots, n-1\}$, to w podciągach $A_1 \dots A_k$ i $A_{k+1} \dots A_n$ rozwiązanie nawiasów musi być optymalne. Wpp. lepsze rozmieszczenie nawiasów w $A_1 \dots A_k$ lub $A_{k+1} \dots A_n$ dałoby lepsze rozmieszczenie nawiasów w całym ciągu.

Rozwiązanie optymalne zawiera rozwiązania optymalne podproblemów (własność optymalnej podstruktury). Podproblemami w naszym problemie są problemy optymalnego rozmieszczenia nawiasów w ciągach $A_i A_{i+1} \dots A_j$ dla $1 \leq i \leq j \leq n$.

Niech $m(i, j)$ oznacza minimalną liczbę mnożeń skalarnych potrzebnych do obliczenia tego iloczynu $A_i \dots A_j$ dla $1 \leq i \leq j \leq n$. Chcemy obliczyć $m(1, n)$.
\begin{flalign*}
	m(i, i) = 0 && &&
\end{flalign*}
Optymalny sposób rozmieszczenia nawiasów jest postaci $(A_i \dots A_k)(A_{k+1} \dots A_j)$ dla pewnego $k \in \{i, \dots, j-1\}$. Stąd mamy zależność:
\begin{flalign*}
	m(i, j) = \min_{i \leq k < j} \left\{ m(i, k) + m(k + 1, j) + p_{i-1} \cdot p_k \cdot p_j \right\} \quad \text{dla } i < j && &&
\end{flalign*}

Niech $s(i, j)$ będzie wartością $k$, dla którego to minimum jest osiągane (czyli $s(i, j)$ wskazuje miejsce ostatniego mnożenia w optymalnym rozwiązaniu podproblemu $A_i \dots A_j$). Czyli dla $1 \leq i < j \leq n$.

Jeśli zaimplementujemy po prostu tę formułę rekurencyjną otrzymamy alg. wykładniczy. Ale jest tylko $n + \binom{n}{2} = \BO(n^2)$ podproblemów, więc rzędu $\BO(n^2)$ liczb $m(i, j)$ i każdą z tych liczb wystarczy obliczyć jeden raz.

\begin{alg}{Algorytm Matrix Chain Order}{alg:matrixchainorder}
	\FunctionNN{MatrixChainOrder}{$p$}
	\Comment{$(p_0, p_1, \dots, p_n)$ - ciąg rozmiarów; $A_i$ - macierz $p_{i-1} \times p_i$}
	\State $n \Assign \Length{p} - 1$
	\For{$i \Assign 1$ \To $n$}
	\State $m(i, i) \Assign 0$
	\EndFor
	\For{$l \Assign 2$ \To $n$} \Comment{$l$ -- liczba macierzy w podproblemie}
	\For{$i \Assign 1$ \To $n - l + 1$} \Comment{$l = 2:\ A_1A_2,\ A_2A_3, \dots, A_{n-1}A_n$}
	\State $j \Assign i + l - 1$ \Comment{$l = 3:\ A_1A_2A_3,\ A_2A_3A_4, \dots, A_{n-2}A_{n-1}A_n$}
	\State $m(i, j) \Assign \infty$ \Comment{$l = n:\ A_1 \dots A_n$}
	\For{$k \Assign i$ \To $j - 1$}
	\State $q \Assign m(i, k) + m(k + 1, j) + p_{i-1} \cdot p_k \cdot p_j$
	\If{$q < m(i, j)$}
	\State $m(i, j) \Assign q$
	\State $s(i, j) \Assign k$
	\EndIf
	\EndFor
	\EndFor
	\EndFor
	\State \Return $m, s$
	\EndFunction
\end{alg}

Mamy 3 pętle (po $l, i, k$) i liczbę iteracji każdej z tych pętli można ograniczyć przez $n$, zatem złożoność czasowa to $\BO(n^3)$.

\begin{example}{Optymalna kolejność mnożenia macierzy dla $n=5$}
$n=5$, $p=(4, 5, 2, 1, 2, 3)$. \\
$A_1 \text{ - } 4 \times 5,\ A_2 \text{ - } 5 \times 2,\ A_3 \text{ - } 2 \times 1,\ A_4 \text{ - } 1 \times 2,\ A_5 \text{ - } 2 \times 3$.

\begin{minipage}{0.45\linewidth}
	$m$
	$\begin{bNiceArray}{c|ccccc}[margin,hvlines]
		i \backslash j & 1 & 2 & 3 & 4 & 5 \\
		1 & 0 & 40 & 30 & 38 & 48 \\
		2 & \times & 0 & 10 & 20 & 31 \\
		3 & \times & \times & 0 & 4 & 12 \\
		4 & \times & \times & \times & 0 & 6 \\
		5 & \times & \times & \times & \times & 0
	\end{bNiceArray}$
\end{minipage}\hfill
\begin{minipage}{0.45\linewidth}
	$s$
	$\begin{bNiceArray}{c|ccccc}[margin,hvlines]
		i \backslash j & 1 & 2 & 3 & 4 & 5 \\
		1 & \times & 1 & 1 & 3 & 3 \\
		2 & \times & \times & 2 & 3 & 3 \\
		3 & \times & \times & \times & 3 & 3 \\
		4 & \times & \times & \times & \times & 4 \\
		5 & \times & \times & \times & \times & \times
	\end{bNiceArray}$
\end{minipage}
\vspace{1em}

\begin{flalign*}
	L=2: \quad & m(1, 2) = 4 \cdot 5 \cdot 2 = 40 && \\
	& m(2, 3) = 5 \cdot 2 \cdot 1 = 10 && \\
	& m(3, 4) = 2 \cdot 1 \cdot 2 = 4 && \\
	& m(4, 5) = 1 \cdot 2 \cdot 3 = 6 && \\[1.5ex]
	%
	L=3: \quad & m(1, 3) = \min \begin{cases}
		m(1, 1) + m(2, 3) + 4 \cdot 5 \cdot 1 \\
		m(1, 2) + m(3, 3) + 4 \cdot 2 \cdot 1
	\end{cases} = \min \begin{cases}
		0 + 10 + 20 \\
		40 + 0 + 8
	\end{cases} = \min \begin{cases} 30 \\ 48 \end{cases} = 30 && \\
	& m(2, 4) = \min \begin{cases}
		m(2, 2) + m(3, 4) + 5 \cdot 2 \cdot 2 \\
		m(2, 3) + m(4, 4) + 5 \cdot 1 \cdot 2
	\end{cases} = \min \begin{cases}
		0 + 4 + 20 \\
		10 + 0 + 10
	\end{cases} = \min \begin{cases} 24 \\ 20 \end{cases} = 20 && \\
	& m(3, 5) = \min \begin{cases}
		m(3, 3) + m(4, 5) + 2 \cdot 1 \cdot 3 \\
		m(3, 4) + m(5, 5) + 2 \cdot 2 \cdot 3
	\end{cases} = \min \begin{cases}
		0 + 6 + 6 \\
		4 + 0 + 12
	\end{cases} = \min \begin{cases} 12 \\ 16 \end{cases} = 12 && \\[1.5ex]
	%
	L=4: \quad & m(1, 4) = \min \begin{cases}
		m(1, 1) + m(2, 4) + 4 \cdot 5 \cdot 2 \\
		m(1, 2) + m(3, 4) + 4 \cdot 2 \cdot 2 \\
		m(1, 3) + m(4, 4) + 4 \cdot 1 \cdot 2
	\end{cases} = \min \begin{cases}
		0 + 20 + 40 \\
		40 + 4 + 16 \\
		30 + 0 + 8
	\end{cases} = \min \begin{cases} 60 \\ 60 \\ 38 \end{cases} = 38 && \\
	& m(2, 5) = \min \begin{cases}
		m(2, 2) + m(3, 5) + 5 \cdot 2 \cdot 3 \\
		m(2, 3) + m(4, 5) + 5 \cdot 1 \cdot 3 \\
		m(2, 4) + m(5, 5) + 5 \cdot 2 \cdot 3
	\end{cases} = \min \begin{cases}
		0 + 12 + 30 \\
		10 + 6 + 15 \\
		20 + 0 + 30
	\end{cases} = \min \begin{cases} 42 \\ 31 \\ 50 \end{cases} = 31 && \\[1.5ex]
	%
	L=5: \quad & m(1, 5) = \min \begin{cases}
		m(1, 1) + m(2, 5) + 4 \cdot 5 \cdot 3 \\
		m(1, 2) + m(3, 5) + 4 \cdot 2 \cdot 5 \\
		m(1, 3) + m(4, 5) + 4 \cdot 1 \cdot 3 \\
		m(1, 4) + m(5, 5) + 4 \cdot 2 \cdot 3
	\end{cases} = \min \begin{cases}
		0 + 31 + 60 \\
		40 + 12 + 24 \\
		30 + 6 + 12 \\
		38 + 0 + 24
	\end{cases} = \min \begin{cases} 91 \\ 76 \\ 48 \\ 62 \end{cases} = 48 &&
\end{flalign*}

Oto algorytm znajdowania optymalnego rozwiązania:
\begin{alg}{Algorytm Matrix Chain Multiply}{alg:matrixchainmultiply}
	\FunctionNN{MatrixChainMultiply}{$A, s, i, j$} \Comment{$A = (A_1, \dots, A_n)$}
	\If{$j = i$}
	\State \Return $A_i$
	\EndIf
	\State $X$ \Assign \Call{MatrixChainMultiply}{$A, s, i, s(i, j)$}
	\State $Y$ \Assign \Call{MatrixChainMultiply}{$A, s, s(i, j) + 1, j$}
	\State \Return \Call{MatrixMultiply}{$X, Y$}
	\EndFunction
\end{alg}

W celu znalezienia rozwiązania wywołujemy: \Call{MatrixChainMultiply}{$A, s, 1, n$}

\begin{flalign*}
	& A_1 \cdot A_2 \cdot A_3 \cdot A_4 \cdot A_5 && \\
	& s(1, 5) = 3 \implies (A_1 \cdot A_2 \cdot A_3)(A_4 \cdot A_5) && \\
	& s(1, 3) = 1 \implies (A_1 \cdot (A_2 \cdot A_3))(A_4 \cdot A_5) \qquad s(4, 5) = 4 \implies \text{bez zmian} && \\
	& s(2, 3) = 2 \implies (A_1 \cdot (A_2 \cdot A_3))(A_4 \cdot A_5) &&
\end{flalign*}
\end{example}

Aby algorytm zadziałał dobrze, problem musi spełniać warunki:
\begin{enumerate}
	\item \textbf{Własność optymalnej podstruktury} -- optymalne rozwiązanie problemu można skonstruować z optymalnych rozwiązań podproblemów.
	\item \textbf{Własność wspólnych podproblemów} -- sytuacja, gdy prosty algorytm rekurencyjny obliczałby rozwiązania tych samych podproblemów wielokrotnie. Liczba podproblemów powinna być ,,mała''.
\end{enumerate}

W klasycznej metodzie programowania dynamicznego rozwiązanie jest budowane od dołu do góry -- od podproblemów najmniejszych do największych. Jest też podejście, w którym rozwiązanie jest budowane od góry do dołu, podobnie jak w prostym algorytmie rekurencyjnym, ale z tą różnicą, że unikamy obliczania tych samych podproblemów wielokrotnie. To podejście nazywa się \textbf{spamiętywaniem} (\textit{memoization}).

\paragraph{Problem mnożenia ciągu macierzy ze spamiętywaniem}
Algorytm oblicza tablice $m$ i $s$ jak poprzednio, ale w inny sposób. Dopóki podproblem nie jest rozwiązany, $m(i, j)$ przechowuje $\infty$. Po rozwiązaniu podproblemu $(i, j)$ umieszczamy w $m(i, j)$ odpowiednią wartość i od tej pory używamy jej bez obliczania.

\begin{alg}{Algorytm Lookup Chain}{alg:lookupchain}
	\FunctionNN{LookupChain}{$p, i, j$}
	\Comment{$p = (p_0, \dots, p_n)$ - ciąg wymiarów; $A_i$ - macierz $p_{i-1} \times p_i$}
	\If{$m(i, j) < \infty$}
	\State \Return $m(i, j), s(i, j)$
	\EndIf
	\If{$i = j$}
	\State $m(i, j) \Assign 0$
	\Else
	\For{$k \Assign i$ \To $j - 1$}
	\State $q \Assign \Call{LookupChain}{p, i, k} + \Call{LookupChain}{p, k + 1, j} + p_{i-1} \cdot p_k \cdot p_j$
	\If{$q < m(i, j)$}
	\State $m(i, j) \Assign q$
	\State $s(i, j) \Assign k$
	\EndIf
	\EndFor
	\EndIf
	\State \Return $(m(i, j), s(i, j))$
	\EndFunction
\end{alg}

\begin{alg}{Algorytm Memoized Matrix Chain}{alg:memoizedmatrixchain}
	\FunctionNN{MemoizedMatrixChain}{$p$}
	\State $n \Assign \Length{p} - 1$
	\For{$i \Assign 1$ \To $n$}
	\For{$j \Assign i$ \To $n$}
	\State $m(i, j) \Assign \infty$
	\EndFor
	\EndFor
	\State \Return \Call{LookupChain}{$p, 1, n$}
	\EndFunction
\end{alg}

\paragraph{Złożoność czasowa:}
\begin{itemize}
	\item Linie 1--4 w \Call{MemoizedMatrixChain}{} zajmują $\BT(n^2)$.
	\item Każdy z $\BO(n^2)$ elementów tablicy $m$ jest modyfikowany raz, później odczytanie wartości $m(i, j)$ zajmuje czas stały.
	\item Modyfikacja $m(i, j)$ zajmuje $\BO(n)$ czasu (pętla w liniach 5--9 w \Call{LookupChain}{}).
\end{itemize}
Ostatecznie złożoność algorytmu wynosi $\BO(n^3)$.


% ════════════════════════════════════════════════════════════════
\subsection{Problem najdłuższego wspólnego podciągu}
% ════════════════════════════════════════════════════════════════

Niech $X = (x_1, \dots, x_m)$ -- ciąg. \\
$Z = (z_1, \dots, z_k)$ jest podciągiem $X$, jeśli istnieje rosnący ciąg indeksów $1 \leq i_1 < i_2 < \dots < i_k \leq m$ takich, że dla każdego $j \in \{1, \dots, k\}$, $z_j = x_{i_j}$.

\begin{example}{Podciąg}
$X = (A, B, C, B, D, A, B)$, $Z = (B, C, D, B)$.
Indeksy: $(2, 3, 5, 7)$.
\end{example}

\paragraph{Problem:}
Dane są ciągi $X = (x_1, \dots, x_m)$ oraz $Y = (y_1, \dots, y_n)$. 
Należy znaleźć najdłuższy wspólny podciąg (LCS -- \textit{Longest Common Subsequence}) $X$ i $Y$, czyli najdłuższy ciąg, który jest podciągiem zarówno $X$, jak i $Y$.

Wszystkich podciągów $X$ jest $2^m$, więc algorytm sprawdzający wszystkie podciągi $X$ działa w czasie wykładniczym.

Niech $X_i = (x_1, \dots, x_i)$ dla $i \in \{1, \dots, m\}$ i dodatkowo $X_0 = ()$ (ciąg pusty). Analogicznie dla $Y$.

\begin{lemma}
	Niech $X = (x_1, \dots, x_m)$, $Y = (y_1, \dots, y_n)$ i niech $Z = (z_1, \dots, z_k)$ będzie LCS $X$ i $Y$.
	\begin{enumerate}
		\item Jeśli $x_m = y_n$, to $z_k = x_m = y_n$ i $Z_{k-1}$ jest LCS $X_{m-1}$ i $Y_{n-1}$.
		\item Jeśli $x_m \neq y_n$ i $z_k \neq x_m$, to $Z$ jest LCS $X_{m-1}$ i $Y$.
		\item Jeśli $x_m \neq y_n$ i $z_k \neq y_n$, to $Z$ jest LCS $X$ i $Y_{n-1}$.
	\end{enumerate}
\end{lemma}

\begin{proof} ~
	\begin{enumerate}
		\item Załóżmy, że $x_m = y_n$. Gdyby $z_k \neq x_m$, to moglibyśmy wydłużyć podciąg $Z$ przez dodanie $z_{k+1} = x_m = y_n$, co daje sprzeczność, bo $(z_1, \dots, z_{k+1})$ jest wspólnym podciągiem $X$ i $Y$ dłuższym niż $Z$. Zatem $z_k = x_m = y_n$ i $Z_{k-1}$ jest wspólnym podciągiem $X_{m-1}$ i $Y_{n-1}$. Gdyby istniał wspólny podciąg $W$ ciągów $X_{m-1}$ i $Y_{n-1}$ długości $> k - 1$, to dodając na koniec $W$ element $x_m = y_n$ otrzymalibyśmy wspólny podciąg $X$ i $Y$ długości $> k - 1 + 1 = k$. Sprzeczność, bo długość $Z$ wynosi $k$.
		\item Załóżmy, że $x_m \neq y_n$ i $z_k \neq x_m$. Wtedy $Z$ jest wspólnym podciągiem ciągów $X_{m-1}$ i $Y$. Gdyby istniał wspólny podciąg $W$ ciągów $X_{m-1}$ i $Y$ o długości $> k$, to $W$ byłby wspólnym podciągiem ciągów $X$ i $Y$ dłuższym niż $Z$. Sprzeczność.
		\item Symetrycznie do punktu 2.
	\end{enumerate}
\end{proof}

Rozwiązanie dla problemu $(X, Y)$ można skonstruować z rozwiązań podproblemów $(X_{m-1}, Y_{n-1})$, $(X_{m-1}, Y)$ i $(X, Y_{n-1})$ -- \textbf{własność optymalnej podstruktury}. 
Podproblem $(X_{m-1}, Y_{n-1})$ jest zawarty w podproblemach $(X_{m-1}, Y)$ i $(X, Y_{n-1})$ -- \textbf{własność wspólnych podproblemów}.

Niech $c(i, j)$ oznacza długość LCS ciągów $X_i$ i $Y_j$.
\begin{flalign*}
	c(i, j) = \begin{cases}
		0 & \text{jeśli } i = 0 \text{ lub } j = 0 \\
		c(i - 1, j - 1) + 1 & \text{jeśli } i, j > 0 \text{ i } x_i = y_j \\
		\max\{c(i - 1, j), c(i, j - 1)\} & \text{jeśli } i, j > 0 \text{ i } x_i \neq y_j
	\end{cases} && &&
\end{flalign*}

Będziemy też używać tablicy $b(i, j)$, aby znaleźć sam podciąg LCS. 
Tablice mają wymiary: $c[0 \dots m, 0 \dots n]$ oraz $b[1 \dots m, 1 \dots n]$. Wartości w tablicy kierunków to $b(i, j) \in \{\nwarrow, \uparrow, \leftarrow\}$.

\begin{alg}{Algorytm obliczający długość LCS}{alg:lcs}
	\FunctionNN{LCS}{$X, Y$}
	\State $m \Assign \Length{X}$
	\State $n \Assign \Length{Y}$
	\For{$i \Assign 0$ \To $m$}
	\State $c(i, 0) \Assign 0$
	\EndFor
	\For{$j \Assign 0$ \To $n$}
	\State $c(0, j) \Assign 0$
	\EndFor
	\For{$i \Assign 1$ \To $m$}
	\For{$j \Assign 1$ \To $n$}
	\If{$x_i = y_j$}
	\State $c(i, j) \Assign c(i - 1, j - 1) + 1$
	\State $b(i, j) \Assign \text{,,}\nwarrow\text{''}$
	\ElsIf{$c(i - 1, j) \geq c(i, j - 1)$}
	\State $c(i, j) \Assign c(i - 1, j)$
	\State $b(i, j) \Assign \text{,,}\uparrow\text{''}$
	\Else
	\State $c(i, j) \Assign c(i, j - 1)$
	\State $b(i, j) \Assign \text{,,}\leftarrow\text{''}$
	\EndIf
	\EndFor
	\EndFor
	\State \Return $(c, b)$
	\EndFunction
\end{alg}

\paragraph{Przykład:}
$X = (A, B, C, B, D, A, B)$, $Y = (B, D, C, A, B, A)$. \\
$m = 7$, $n = 6$.

\begin{table}[H]
	\centering
	\renewcommand{\arraystretch}{1.2}
	\begin{NiceTabular}{cc|c|cccccc}[margin,hvlines]
		& $j$ & 0 & 1 & 2 & 3 & 4 & 5 & 6 \\
		$i$ & & $y_j$ & $B$ & $D$ & $C$ & $A$ & $B$ & $A$ \\
		\hline
		0 & $x_i$ & 0 & 0 & 0 & 0 & 0 & 0 & 0 \\
		1 & $A$   & 0 & $\uparrow 0$ & $\uparrow 0$ & $\uparrow 0$ & $\nwarrow 1$ & $\leftarrow 1$ & $\nwarrow 1$ \\
		2 & $B$   & 0 & $\nwarrow 1$ & $\leftarrow 1$ & $\leftarrow 1$ & $\uparrow 1$ & $\nwarrow 2$ & $\leftarrow 2$ \\
		3 & $C$   & 0 & $\uparrow 1$ & $\uparrow 1$ & $\nwarrow 2$ & $\leftarrow 2$ & $\uparrow 2$ & $\uparrow 2$ \\
		4 & $B$   & 0 & $\nwarrow 1$ & $\uparrow 1$ & $\uparrow 2$ & $\uparrow 2$ & $\nwarrow 3$ & $\leftarrow 3$ \\
		5 & $D$   & 0 & $\uparrow 1$ & $\nwarrow 2$ & $\uparrow 2$ & $\uparrow 2$ & $\uparrow 3$ & $\uparrow 3$ \\
		6 & $A$   & 0 & $\uparrow 1$ & $\uparrow 2$ & $\uparrow 2$ & $\nwarrow 3$ & $\uparrow 3$ & $\nwarrow 4$ \\
		7 & $B$   & 0 & $\nwarrow 1$ & $\uparrow 2$ & $\uparrow 2$ & $\uparrow 3$ & $\nwarrow 4$ & $\uparrow 4$ \\
	\end{NiceTabular}
\end{table}

\begin{alg}{Wypisywanie najdłuższego wspólnego podciągu}{alg:printlcs}
	\FunctionNN{PrintLCS}{$b, X, Y, i, j$}
	\If{$i = 0$ \textbf{lub} $j = 0$} 
	\State \Return $\varepsilon$ \Comment{ciąg pusty}
	\EndIf
	\If{$b(i, j) = \text{,,\nwarrow''}$}
	\State \Call{PrintLCS}{$b, X, Y, i - 1, j - 1$}
	\State \textbf{print} $x_i$
	\ElsIf{$b(i, j) = \text{,,\uparrow''}$}
	\State \Call{PrintLCS}{$b, X, Y, i - 1, j$}
	\Else
	\State \Call{PrintLCS}{$b, X, Y, i, j - 1$}
	\EndIf
	\EndFunction
\end{alg}

Żeby wypisać rozwiązanie, wywołujemy \Call{PrintLCS}{$b, X, Y, m, n$}. \\
Dla powyższego przykładu otrzymamy: $(B, C, B, A)$.

Złożoność czasowa tego algorytmu wynosi $\BT(m \cdot n)$.